\subsection{Executing the UBS notice and conversion terms}
\label{sec:ubs-terms}

The UBS example can now begin with a particular issue rather than a generic
conversion rule. Annex A of the retained published terms concerns the
SGD 500 million notes issued on 24 June 2024, with ISIN CH1357852636.
The same PDF contains the later SGD 175 million reopening in Annex B.
The initial conversion price for the June issue is SGD 37.77, and holdings
must respect the SGD 250,000 denomination and increment. Keeping the annex
explicit prevents an internal capital loan, another currency issue or the
reopening's issue date from supplying the wrong premise
\citep[Annex A, definitions and condition~2]{ubsat12024}.

These terms describe direct subordinated obligations that can convert into
shares. The source-dependent product reading therefore supplies five previously
missing premises: debt form, qualifying contingent loss absorption, absence of
the plain-debt or deposit exclusion, and absence of either a related-product
wrapper or a loss-absorption fund. They yield a conditional conclusion that the
direct notes fall within the selected HKMA product scope. The reading remains
an interpretation of the retained edition. It does not establish current legal
force, the purchaser's PI status or the bank's capacity. The combined assessment
continues to carry those questions through all fourteen bank obligations.

\paragraph{A breach creates a notice question.}
The capital test uses the figures attached to a particular publication.
Let $C$ be its CET1 capital, $H$ the defined Higher-Trigger Amount, and $R>0$
its risk-weighted assets, all expressed in the same monetary unit. The
Higher-Trigger Amount represents capital capable of absorbing losses under
the higher-trigger provisions; omitting it can change the result. The
published ratio and breach test are
\begin{equation}
  r=\frac{C+H}{R},\qquad
  r<0.07\quad\Longleftrightarrow\quad 100(C+H)<7R.
  \label{eq:ubs-trigger-ratio}
\end{equation}
The positive denominator permits multiplication without reversing the
inequality. With $C=6$, $H=1$ and $R=100$, the ratio is exactly seven percent,
so this strict test does not breach. With $H=0$, it does. Condition 7(d)
uses the publication's figures even if they are later restated. The software
binds the calculation to that retained publication snapshot
\citep[definition of Trigger CET1 Ratio; conditions~7(b), 7(d)]{ubsat12024}.

A breach is not itself the defined Trigger Event. Condition 7(a) ties that
event to the issuer giving a Trigger Event Notice in accordance with the
terms. An ordinary publication normally starts a five-business-day notice
period; an extraordinary publication calls for notice on the publication
date. The ordinary branch also has a restoration exception involving a UBS
request, FINMA's written agreement, restoration or imminent restoration above
the threshold, their assessment of adequacy, and a strict timing condition.
The agreement must precede the earlier of the ordinary notice date and the
original fifth-business-day determination date. An unknown earlier notice
history or unresolved ordering within one day therefore leaves the exception
qualified. A bare restoration percentage cannot establish it.

Outstanding higher-trigger instruments can postpone both the notice and the
conversion. The program retains the full supplied inventory and selects the
latest relevant conversion date; one observed higher-trigger notice cannot
stand for all of them. Notice publication is a further factual question.
Condition 14 uses the first qualifying SIX publication while the notes are
listed, or delivery to the intermediary after delisting. The implemented
publication path does not infer a notice from its deadline; alternative
exchange-rule notice methods still require their own evidence and treatment
\citep[conditions~7(a)--(b), 14]{ubsat12024}.

\paragraph{Different events use different clocks.}
For the publication definitions, the Business Day definition refers to Zurich.
For the other relevant calculations it requires Singapore and Zurich.
The calendar input therefore records a covered interval and the days on which
the specified banks and foreign-exchange markets are open in each location.
An ordinary weekday or a public-holiday list alone does not establish that
condition. The engineering examples use explicitly hypothetical calendars.

A Viability Event has its own FINMA and, where relevant, public-support
premises. Its notice must be given within three days, while conversion must
occur within twenty Business Days following the event. The latter period
starts from the event, not from the notice. In a hypothetical calendar with
both markets open on every weekday, a Friday event on 4 September 2026 has a
Monday notice deadline on 7 September and a conversion deadline on 2 October.
Starting the twenty-day count from that Monday notice would instead produce
5 October. The program rejects that extra time. A qualifying change in law
can stop the viability provisions under Condition 7(e), but they cease from
the specified actual notice date, not merely the date of the joint
determination. The capital-ratio trigger remains a separate branch
\citep[Business Day definition; conditions~7(c), 7(e)]{ubsat12024}.

\paragraph{Shares, prices and delivery need separate evidence.}
Once the appropriate notice and price premises are supplied, the calculation
aggregates the principal held by one holder before discarding any fraction
of a share. For aggregate principal $P$ and conversion price $K>0$ in the
same currency, the number of shares is
\begin{equation}
  N=\left\lfloor\frac{P}{K}\right\rfloor,
  \qquad NK\leq P<(N+1)K.
  \label{eq:ubs-whole-shares}
\end{equation}
The inequality checks the rounding independently. Rounding each lot first
can discard more than one holder's terms permit. The discarded fraction
does not produce cash compensation. Nor does the resulting share count
establish a market value. Actual share creation, receipt by the settlement
depository, cancellation of accrued unpaid interest and subsequent delivery
to the holder remain distinct events
\citep[conditions~8(a)--(c), 8(g)--(j)]{ubsat12024}.

The price is adjusted through the notice date, including that date. A later
corporate action cannot silently replace the price applicable to the notice.
Implemented arithmetic covers declared share splits and bonus issues,
discounted rights, exact-cent adjustments, the par-value floor and a bounded
carry-forward case. Other calculations cover the stated successor-share price
ratio and the distribution of supplied net offer proceeds. Their market
prices, exchange rates, relevant dealing days and contractual determinations
remain explicit premises. Ambiguous rounding, overlapping adjustments,
unavailable action history, amended terms and issuer substitution qualify
the result rather than select a convenient convention.

The retained source itself supplies a useful challenge. Condition 8(d)(i)(C)
prints the extraordinary-distribution factor as $(A-B)/B$, where $A$ is the
specified market price and $B$ is the distribution per share. The rendered
PDF confirms the denominator. At $A=100$ and $B=10$, that expression gives
$9$, whereas $(A-B)/A$ would give $0.9$. The program preserves the printed
expression and its literal result, records the unresolved intended adjustment,
and withholds an operational price for that branch. It does not correct the
contract by appealing to the more familiar formula
\citep[condition~8(d)(i)(C), PDF p.~30]{ubsat12024}.

The checks exercise these distinctions without using human expected legal
answers. Independent mathematical relations test the capital comparison,
restoration predicates, date ordering, product scope and PI restriction;
native RuleIR/Java and Catala executions test the corresponding programs.
Calendar selection and whole-share inequalities separately challenge the
Python calculations. Source substitutions, capital restatements, missing
notices and changed scenarios must invalidate or qualify the affected result.
The PI and non-PI examples change declared client premises while holding the
issue fixed. They distinguish the selected selling condition without turning
either example into approval by an actual private bank. Incorporated agency
agreements, actual notices and market histories, applicable current law and
private account evidence remain required for a real transaction.
